\documentclass{beamer}

\mode<presentation>
{
\usetheme{Darmstadt}
\setbeamercovered{transparent}
}

\usepackage[english]{babel}
\usepackage[latin1]{inputenc}
\usepackage{mathptmx}
\usepackage[scaled=.90]{helvet}
\usepackage{courier}
\usepackage[T1]{fontenc}
\usepackage{pgf}
\usepackage{tikz}
\usepackage{common}
\usepackage[normalem]{ulem}

\newcommand{\hcancel}[1]{
    \tikz[baseline=(tocancel.base)]{
        \node[inner sep=0pt,outer sep=0pt] (tocancel) {#1};
        \draw[red] (tocancel.south west) -- (tocancel.north east);
    }
}

\pgfdeclareimage[width=11cm,height=6cm]{any_questions}{any_questions}

\title{Computing All Squares in Compressed Texts}

\author{Lesha Khvorost\inst{\mbox{\textdagger}}}

\institute
{
\inst{\mbox{\textdagger}}
Institute of Mathematics and Computer Science\\
Ural Federal University}

\date{26th September RuFiDiM 2012}
\logo{\includegraphics[height=0.8cm]{logo.jpg}}
\subject{Talks}

\AtBeginSubsection[]
{
\begin{frame}<beamer>
\frametitle{Outline}
\tableofcontents[currentsection,currentsubsection]
\end{frame}
}

\begin{document}

\begin{frame}
\titlepage
\end{frame}

\section{Introduction}

\begin{frame}
\frametitle{}

\begin{block}{}
String algorithms are deeply explored.
\end{block}

\pause
Two research directions are popular now:

\begin{itemize}
  \item To improve classic string algorithms to process huge inputs (\textbf{I/O efficient algorithms});
\end{itemize}

\pause

\begin{block}{}
ATLAS project at CERN produces 70Tb/s of valuable information.
\end{block}

\pause

\begin{itemize}
  \item To change data representation from the native to a more compact one (\textbf{algorithms on compressed representations}).
\end{itemize}
\end{frame}

\begin{frame}
\frametitle{Features of algorithms on compressed representations}

\begin{itemize}
  \item We save storage space and measure the complexity of the algorithm in terms of compressed input;
  \pause
  \item By default we work with preprocessed data.
\end{itemize}

\end{frame}

\section{SLPs Background}

\begin{frame}
\frametitle{}

\begin{definition}
    A \textbf{straight-line program} (SLP) $\slp{S}$ is a sequence of assignments of the form:
    
    \begin{center}
        $\slp{S}_1 = expr_1,\ \slp{S}_2 = expr_2, \dots, \slp{S}_n = expr_n,$
    \end{center}
    where $\slp{S}_i$ are \textbf{rules} and $expr_i$ are expressions of the form:
    \begin{itemize}
        \item $expr_i \in \Sigma$ (\textbf{terminal} rules), or
        \item $expr_i = \slp{S}_\ell\cdot \slp{S}_r \ (\ell, r < i)$
        (\textbf{nonterminal} rules).
    \end{itemize}
\end{definition}

\pause

\begin{example}[the SLP that derives the 6th Fibonacci word]
    \begin{center}
    $\slp{F}_0{\to} b,\ \slp{F}_1{\to} a,\ \slp{F}_2{\to} \slp{F}_1\cdot\slp{F}_0,\ \slp{F}_3{\to} \slp{F}_2\cdot\slp{F}_1,$ 
    
    $\slp{F}_4{\to} \slp{F}_3\cdot\slp{F}_2,\ \ \slp{F}_5{\to} \slp{F}_4\cdot\slp{F}_3,\ \ \slp{F}_6{\to} \slp{F}_5\cdot\slp{F}_4.$
    \end{center}
\end{example}

\end{frame}

\begin{frame}
\frametitle{}

\begin{example}[the parse of the SLP that derives the 6th Fibonacci word]
    \begin{center}
        \FibonacciWordSLP
    \end{center}
\end{example}

\end{frame}

\begin{frame}
\frametitle{Tools}

\begin{itemize}
  \item Collecting some information from children (\textbf{Pattern Matching}, \textbf{Longest Common Substring});
\end{itemize}

\pause 

\begin{center}
    \collectingTool
\end{center}

\pause

\begin{alertblock}{}
Sometimes count operations are extremely expensive (\textbf{Hamming Distance}, \textbf{Levenstein Distance}).
\end{alertblock}

\pause

\begin{block}{}
Are there other tools for SLPs processing?
\end{block}

\end{frame}

\section{CAS Problem}

\begin{frame}
\frametitle{}

\begin{definition}
A string is called a square if it can be obtained by concatenating two copies of some string.\newline
For example, \emph{abc} $\cdot$ \emph{abc}. We say that \emph{abc} is the root of the square.
\end{definition}

\pause

\begin{block}{CAS Problem}
\textsc{Input:} an SLP $\slp{S}$ that derives the text $S$;

\textsc{Output:} the $S$-table that holds information about all squares in $S$.
\end{block}

\pause

Main difficulties:
\begin{itemize}
  \item The text may contains exponentially many squares (in length of a compressed representation) 
  \newline ($a^n$ contains $O(n^2)$ squares and has a $\log n$ representation);
  \pause
  \item It is hard to collect information from children.
\end{itemize}

\end{frame}

\begin{frame}
\begin{definition}
A string S is called a palindrome if $S = S^R$ (for example \emph{ababa}).
\end{definition}

\pause

\begin{block}{Computing All Palindromes Problem}
\textsc{Input:} an SLP $\slp{S}$ that derives the text $S$;

\textsc{Output:} a data structure that holds information about all palindromes in $S$.
\end{block}

\pause

\begin{block}{Theorem [W. Matsubara, S. Inenaga et al] (2008)}
Computing all palindromes problem can be solved in $O(|\slp{S}|^4)$ time with $O(|\slp{S}|^2)$ space.
\end{block}

\end{frame}

\begin{frame}
\frametitle{Main result}

\begin{block}{Theorem}
There is an algorithm that solves CAS using $O(|\slp{S}|^4 \cdot \log^2{|S|})$ time and $O(|\slp{S}| \cdot \max(|\slp{S}|, \log{|S|}))$ space.
\end{block}

\end{frame}

\begin{frame}
\frametitle{S-table}

\begin{block}{}
The \emph{S-table} is a rectangular table $S(\slp{S})$ of size $(\lfloor\log |S|\rfloor+1) \times
(|\slp{S}|+1)$ that stores families of squares. 

\pause

There are three types of families: 
\begin{itemize}
  \item Simple family, parameterized by $\{|x|, c_\ell, c_r\}$;
  \pause
  \item Two types of complex families.
\end{itemize}
\end{block}

\pause 

\begin{example}[Simple family of squares]
$S = bbabbabb$, family: $\{3, 3, 5\}$, squares: $\{(bba)^2, (bab)^2, (abb)^2 \}$.
\end{example}

\end{frame}

\begin{frame}
\frametitle{Problems solvable by usage of $S$-table}

For a given $\slp{S}$ that derives a text S
\begin{itemize}
  \item to check whether $S$ is a square or not ($O(1)$);
  \pause
  \item to find information about all squares of fixed length ($O(|\slp{S}|)$);
  \pause
  \item to check whether on not a text $S$ is square-free ($O(|\slp{S}| \log |S|)$);
  \pause
  \item to compute total number of squares that are contained in $S$ ($O(|\slp{S}| \log |S|)$);
  \pause
  \item to find a maximal by length square that occurs in $S$ ($O(|\slp{S}| \log |S|)$);
  \pause
  \item to find information about all squares that starts from a fixed position~$i$ ($O(|\slp{S}| \log |S|)$).
\end{itemize}

\pause

\begin{alertblock}{Open problem}
For a given $\slp{S}$ that derives a text $S$ to construct an SLP that derives concatenation of all squares that contain in $S$.
\end{alertblock}

\end{frame}

\begin{frame}
\frametitle{Detecting squares}

\begin{center}
    \squaresDetecting
\end{center}
\end{frame}

\section*{Conclusion}

\begin{frame}
\frametitle{Conclusion}

\begin{itemize}
  \item<1-> \sout{We present the algorithm of 6th power that solves CAS Problem}
  \item<2-> We present the new way of SLPs processing
\end{itemize}

\pause

\begin{alertblock}{Open problem}<3->
Is it possible to improve the algorithm's complexity? 
\end{alertblock}
\end{frame}

\begin{frame}
\begin{center}
    \begin{tikzpicture}
        \draw (0,0) node {\pgfuseimage{any_questions}};
    \end{tikzpicture}
\end{center}

\end{frame}

\begin{frame}
\frametitle{Serching squares}

\begin{center}
    \squaresAlgorithm
\end{center}

\end{frame}

\begin{frame}
\frametitle{Palindromes location}

Let $\slp{S}$ be equal to $\slp{S}_\ell \cdot \slp{S}_r$

\palindromesLocation

\begin{alertblock}{}<5->
How we can use infromation about squares from children?
\end{alertblock}

\end{frame}

\end{document}

